\chapter{System Requirements and Design Drivers}

 
\section{Overview}

The ND280 Magnetic Field Monitoring System has been designed to satisfy a combination of scientific, functional, operational and engineering requirements.

The requirements defined in this chapter provide the reference framework for the hardware architecture, NodeOS firmware, ND280 Monitor application, validation activities and future system extensions.

Each requirement is assigned a unique identifier to support traceability throughout the development and qualification process.

\section{Functional Requirements}

\begin{longtable}{p{18mm}p{106mm}p{23mm}}
	\toprule
	ID & Requirement & Status \\
	\midrule
	
	REQ-001 &
	The system shall continuously acquire the three magnetic-field components $B_x$, $B_y$ and $B_z$. &
	Verified \\
	
	REQ-002 &
	The system shall acquire the temperature measured by each magnetic sensor. &
	Verified \\
	
	REQ-003 &
	Each acquisition node shall provide a unique hardware identifier. &
	Verified \\
	
	REQ-004 &
	Calibration parameters shall be stored in non-volatile memory and shall survive power cycling. &
	Verified \\
	
	REQ-005 &
	Each acquisition node shall automatically recover from a firmware stall through hardware watchdog supervision. &
	Verified \\
	
	REQ-006 &
	Each acquisition node shall provide diagnostic information including acquisition statistics, communication errors and reset cause. &
	Verified \\
	
	REQ-007 &
	The monitoring software shall provide real-time visualization of magnetic-field and temperature measurements. &
	Verified \\
	
	REQ-008 &
	The monitoring software shall provide persistent data logging for long-duration acquisitions. &
	Verified \\
	
	REQ-009 &
	The command interface shall reject malformed or unrelated serial traffic without disrupting the acquisition process. &
	Verified \\
	
	REQ-010 &
	The magnetic-field measurement range shall provide adequate headroom with respect to the nominal ND280 magnetic field of approximately \SI{0.2}{T}. &
	Verified \\
	
	REQ-011 &
	The local acquisition architecture shall support up to eight TMAG5273A2 sensors connected to a single intelligent node. &
	Rev.A design \\
	
	REQ-012 &
	The system architecture shall support multiple acquisition nodes connected through a common inter-board communication network. &
	In development \\
	
	REQ-013 &
	The distributed architecture shall support up to 32 acquisition boards, corresponding to a maximum of 256 magnetic sensors. &
	Design target \\
	
	REQ-014 &
	The same acquisition-board hardware shall be capable of operating either as a NODE or as a COORDINATOR. &
	Rev.A design \\
	
	REQ-015 &
	The NODE or COORDINATOR operating role shall be selected through a dedicated hardware configuration input. &
	Rev.A design \\
	
	\bottomrule
\end{longtable}

\section{Reliability Requirements}

\begin{longtable}{p{18mm}p{106mm}p{23mm}}
	\toprule
	ID & Requirement & Status \\
	\midrule
	
	REL-001 &
	The acquisition process shall continue without operator intervention during normal long-duration operation. &
	Verified \\
	
	REL-002 &
	A communication error shall not corrupt the internal measurement state of the acquisition node. &
	Verified \\
	
	REL-003 &
	Invalid EEPROM contents shall be detected and replaced by safe default calibration values. &
	Verified \\
	
	REL-004 &
	The firmware shall record the cause of the most recent microcontroller reset. &
	Verified \\
	
	REL-005 &
	The external I\textsuperscript{2}C interface of each board shall remain disabled during microcontroller reset and early boot. &
	Rev.A design \\
	
	REL-006 &
	The inter-board communication architecture shall electrically segment successive sections of the I\textsuperscript{2}C network. &
	Rev.A design \\
	
	\bottomrule
\end{longtable}

\section{Engineering Design Drivers}

The system architecture has also been influenced by a number of engineering design drivers that are not purely functional requirements but are considered essential for maintainability and future system evolution.

\begin{itemize}[leftmargin=*]
	
	\item \textbf{Modularity.}
	Hardware abstraction, device drivers, acquisition services and application logic shall remain separated.
	
	\item \textbf{Firmware reuse.}
	A new hardware revision should not require a complete redesign of NodeOS.
	
	\item \textbf{Deterministic embedded behaviour.}
	Dynamic memory allocation is avoided in the embedded firmware.
	
	\item \textbf{Hardware-defined identity.}
	The Node ID is determined by hardware configuration rather than by software or EEPROM contents.
	
	\item \textbf{Hardware-defined operating role.}
	The NODE/COORDINATOR role is determined by a dedicated hardware jumper.
	
	\item \textbf{Local data processing.}
	Each intelligent node performs sensor acquisition, averaging, calibration and diagnostics locally before data are transferred to higher-level controllers.
	
	\item \textbf{Independent subsystem validation.}
	Each hardware and software subsystem should be testable independently.
	
	\item \textbf{Incremental qualification.}
	New functionality is introduced only after the previous development stage has been validated.
	
	\item \textbf{Single-board architecture.}
	Whenever practical, NODE and COORDINATOR operation should use the same PCB and bill of materials.
	
\end{itemize}


\section{Calibration and Metrological Requirements}

The ND280 Magnetic Field Monitoring System is not intended to replace a dedicated high-precision magnetic-field mapping instrument. Its primary purpose is distributed, long-term monitoring of the magnetic field and the detection of spatial or temporal variations relevant to detector operation. Nevertheless, the usefulness of such a monitor depends on the measurements being sufficiently stable, reproducible and comparable between sensors and acquisition nodes.

A recent calibration study of a dedicated three-dimensional Hall sensor card developed at CERN illustrates the principal error sources that can affect Hall-based vector measurements, including offset, temperature dependence, non-linearity, cross-axis effects and mechanical or angular misalignment. The same study emphasizes that a dedicated calibration procedure is required when high measurement accuracy is needed \cite{MessnerHallCalibration}. The absolute performance target of that system is substantially more demanding than required for the ND280 monitoring application; however, the underlying metrological methodology is directly relevant to the present project.

The following requirements are therefore adopted for the calibration and characterization programme of the ND280 monitoring nodes:

\begin{longtable}{p{18mm}p{106mm}p{23mm}}
\toprule
ID & Requirement & Status \\
\midrule

CAL-001 &
Each magnetic sensor shall have an individually identifiable calibration data set. &
Design target \\

CAL-002 &
The calibration model shall support correction of per-axis offset and sensitivity and shall be extensible to cross-axis and alignment corrections if required by characterization data. &
Design target \\

CAL-003 &
Calibration coefficients shall be stored in non-volatile memory together with a calibration format or version identifier. &
Partially implemented \\

CAL-004 &
The calibration procedure shall characterize the temperature dependence of the magnetic-field measurement over the expected operating temperature range. &
Planned \\

CAL-005 &
The calibration campaign shall evaluate sensor-to-sensor and board-to-board reproducibility using multiple independently assembled units. &
Planned \\

CAL-006 &
The calibration model shall be validated over the magnetic-field range relevant to ND280 operation and shall provide adequate margin around the nominal field of approximately \SI{0.2}{T}. &
Planned \\

CAL-007 &
The system shall preserve the raw or minimally processed sensor information required to repeat calibration studies and to evaluate alternative correction models offline. &
Design target \\

CAL-008 &
The calibration procedure and its resulting coefficients shall be traceable through documented calibration conditions, sensor identity, calibration date and validation results. &
Planned \\

\bottomrule
\end{longtable}

The calibration model will intentionally be kept no more complex than justified by the characterization results. The current offset-and-gain representation therefore remains a valid first-level model. A full three-dimensional correction matrix and explicit temperature terms will be introduced only if measurements demonstrate that they provide a significant improvement in the operating range of the monitoring system.

\section{Scalability}

The system has been intentionally developed through successive levels of complexity:

\begin{enumerate}[leftmargin=*]
	
	\item Reference Single-Sensor Node;
	
	\item Eight-Sensor Multiplexing Node;
	
	\item Distributed Multi-Board Acquisition System;
	
	\item Complete detector monitoring infrastructure.
	
\end{enumerate}

Each validated development stage provides the reference platform for the following stage.

The Reference Single-Sensor Node validates the complete acquisition chain and the NodeOS architecture.

The Eight-Sensor Multiplexing Node extends local acquisition capability by introducing the TCA9548A while preserving the existing firmware architecture.

The distributed system introduces inter-board communication and coordinator functionality without modifying the local sensor-acquisition concept.

\section{Verification Strategy}

Each requirement shall be associated with one or more verification activities.

Verification methods include:

\begin{itemize}[leftmargin=*]
	
	\item schematic and PCB design review;
	
	\item firmware unit and integration tests;
	
	\item communication protocol tests;
	
	\item watchdog and reset-recovery tests;
	
	\item EEPROM persistence tests;
	
	\item sensor acquisition tests;
	
	\item long-duration stability tests;
	
	\item inter-board communication tests;
	
	\item final detector-level characterization.
	
\end{itemize}

A Requirements Traceability Matrix will be maintained as part of the validation documentation and will associate each requirement with the corresponding implementation and verification evidence.

\section{Summary}

The requirements defined in this chapter establish the functional and engineering framework of the ND280 Magnetic Field Monitoring System.

The design prioritizes measurement continuity, robustness, modularity, scalability and long-term maintainability. These requirements directly motivate the system architecture described in the following chapter.